\chapter{典型裂项题型}

\section{反正切裂项}

基本公式：$$ \arctan X - \arctan Y = \arctan \frac{X-Y}{1+XY} $$$$ \operatorname{arccot} X - \operatorname{arccot} Y = \operatorname{arccot} \frac{XY+1}{Y-X} $$

\begin{example}[$\sum \arctan \frac{2n^2-4}{n^4+n^2+3}$]
    \[\sum_{n=1}^{\infty} \arctan \frac{2n^2-4}{n^4+n^2+3}\]
\end{example}
\begin{solution}
取
\[
X = \frac{n-1}{(n-1)^2+1},\qquad Y = \frac{n+1}{(n+1)^2+1}.
\]
先分别通分 \(X-Y\) 与 \(1+XY\)。两者分子中都出现公共因子
\[
n^4+4=(n^2-2n+2)(n^2+2n+2),
\]
约去这个因子后，差角公式中的分式与题中的分式一致。

由基本公式 \(\arctan X - \arctan Y = \arctan \frac{X-Y}{1+XY}\) 可得，代入上述 \(X,Y\) 后，原级数化为：
\[
\sum_{n=1}^{\infty} \left( \arctan \frac{n-1}{(n-1)^2+1} - \arctan \frac{n+1}{(n+1)^2+1} \right) = \arctan 0 - \lim_{n\to\infty} \arctan \frac{n+1}{(n+1)^2+1} = 0 - 0 = 0
\]

\end{solution}

\begin{example}[$\operatorname{arccot} \frac{a_{n+1}}{a_n} - \operatorname{arccot} \frac{a_n}{a_{n-1}}$]
    递推式 \( a_{n+1} = 4a_n - a_{n-1} \)，首项 \( a_1=2, a_2=8 \)。求
    \[\sum_{n=2}^{\infty} \operatorname{arccot} a_n^2\]
\end{example}
\begin{solution}
令
\[
X = \frac{a_{n+1}}{a_n},\qquad Y = \frac{a_n}{a_{n-1}}.
\]
要把反余切差角公式落到题中这一项，关键是验证
\[
\frac{XY+1}{Y-X}=a_n^2.
\]
其中分子由递推式化简，分母由不变量
\[
a_{n+1}a_{n-1}-a_n^2=-4
\]
化简，于是右端确实等于 \(a_n^2\)，裂项结构由此确定。

代入 \( X,Y \) 后，\(\frac{XY+1}{Y-X} = a_n^2\)，因此原级数可写成 \(\sum (\operatorname{arccot} X - \operatorname{arccot} Y)\)。求和得：
\[
\operatorname{arccot} \frac{a_2}{a_1} - \lim_{n\to\infty} \operatorname{arccot} \frac{a_{n+1}}{a_n} = \operatorname{arccot} 4 - 0 = \frac{\pi}{2} - \arctan 4
\]

\end{solution}

\begin{example}[$\sum_{n=-\infty}^{+\infty} \arctan\left(\frac{1}{\sqrt{6}F_n + 1}\right)$]
    设数列\(F_n\)满足\(F_1 = F_2 = 1\)且对任意整数\(n\)都有\(F_n = F_{n-1} + F_{n-2}\)，证明：
    $$\sum_{n=-\infty}^{+\infty} \arctan\left(\frac{1}{\sqrt{6}F_n + 1}\right) = \sum_{n=-\infty}^{+\infty} \arctan\left(\frac{1}{\sqrt{6}F_n - 1}\right) = \frac{\pi}{2}$$
\end{example}
\begin{solution}
取
\[
X=\sqrt{\frac{3}{2}}F_{2n},\qquad Y=\sqrt{\frac{3}{2}}F_{2n-2}.
\]
再利用
\[
\arctan X-\arctan Y=\arctan\frac{X-Y}{1+XY}
\]
以及卡西尼恒等式
\[
F_{2n}F_{2n-2}=F_{2n-1}^2-1,
\]
可把差角公式中的分母整理成 \(1+XY\) 的形状，因此每一项都写成相邻两项反正切之差。

令 \( X_n = \sqrt{\frac{3}{2}}F_{2n} \)，则每一项等于 \(\arctan X_n - \arctan X_{n-1}\)。对对称部分和求极限，得第一和为 \(\frac{\pi}{2}\)。把上面的 \(X_n\) 改为 \(-X_n\) 时，差角公式中的 \(+1\) 相应变为 \(-1\)，同样得到一组首尾相消的反正切差，因此第二和也为 \(\frac{\pi}{2}\)。
\end{solution}

\begin{example}[$\frac{\sinh\pi}{\cosh(2n\pi)-\cosh\pi}$]
\(\displaystyle\sum_{n=1}^{\infty}\frac{\sinh\pi}{\cosh(2n\pi)-\cosh\pi}\)，这里\(\sinh x=\frac{e^x-e^{-x}}{2}\), \(\cosh x=\frac{e^x+e^{-x}}{2}\)
\end{example}
\begin{solution}
令
\[
A=\left(n-\frac12\right)\pi,\qquad B=\left(n+\frac12\right)\pi.
\]
由恒等式
\[
\coth A-\coth B=\frac{\sinh(B-A)}{\sinh A\,\sinh B}
\]
并配合双曲函数的积化和差公式，可把原项改写成相邻两项 \(\coth\) 的差。

原项等于 \(\frac{1}{2}(\coth\frac{\pi}{2} - \coth(n+\frac{1}{2})\pi)\)。求和得：
\[
\frac{1}{2}\left(\coth\frac{\pi}{2} - \lim_{n\to\infty}\coth\left(n+\frac{1}{2}\right)\pi\right) = \frac{1}{2}\left(\coth\frac{\pi}{2} - 1\right) = \frac{1}{e^\pi - 1}
\]

\end{solution}

\begin{example}[$\sum_{n=2}^{\infty} \operatorname{arccot} (3a_n^2) $]
设数列满足 \(a_1=1, a_2=3, a_{n+1} = 3a_n - a_{n-1}\)。计算级数：$$ \sum_{n=2}^{\infty} \operatorname{arccot} (3a_n^2) $$
\end{example}
\begin{solution}
令
\[
X=\frac{a_{n+1}}{a_n},\qquad Y=\frac{a_n}{a_{n-1}}.
\]
由递推式 \(a_{n+1}=3a_n-a_{n-1}\) 以及不变量
\[
a_{n+1}a_{n-1}-a_n^2=-1
\]
可化简
\[
\frac{XY+1}{Y-X}=3a_n^2,
\]
从而得到反余切裂项。

求和得：
\[
\operatorname{arccot} 3 - \lim_{n\to\infty} \operatorname{arccot} (3a_{n+1}/a_n) = \operatorname{arccot} 3 - 0 = \operatorname{arccot} 3
\]

\end{solution}

\begin{example}[$ \sum_{n=1}^{\infty} \arctan \left( \frac{n^2+n-1}{n^4+2n^3+4n^2+3n+2} \right) $]
    $$ \sum_{n=1}^{\infty} \arctan \left( \frac{n^2+n-1}{n^4+2n^3+4n^2+3n+2} \right) $$
\end{example}
\begin{solution}
取
\[
X=\frac{n}{n^2+1},\qquad Y=\frac{n+1}{n^2+2n+2}.
\]
把 \(X,Y\) 代入差角公式，先化分子 \(X-Y\)，再化分母 \(1+XY\)，二者合并后得到题中的分式，因此原级数可以写成反正切差。

因此原级数化为 \(\sum (\arctan X - \arctan Y)\)，从而
\[
\arctan 1 - \lim_{n\to\infty} \arctan Y = \frac{\pi}{4}.
\]
\end{solution}

\section{多项式与根式型}

\begin{example}[$\displaystyle \sum_{n=1}^{99} \sqrt{1 + \frac{1}{n^2} + \frac{1}{(n+1)^2}}$]
    $$\displaystyle \sum_{n=1}^{99} \sqrt{1 + \frac{1}{n^2} + \frac{1}{(n+1)^2}}$$
\end{example}
\begin{solution}
利用完全平方公式 \((a+b-c)^2\)。令 \( a=1, b=1/n, c=1/(n+1) \)，交叉项 \( 2(ab-ac-bc) \) 恰好抵消，因此 \( 1 + 1/n^2 + 1/(n+1)^2 = (1 + 1/n - 1/(n+1))^2 \)。

开根号得 \( \left|1 + \frac{1}{n} - \frac{1}{n+1}\right| = 1 + \frac{1}{n} - \frac{1}{n+1} \)，因此
\[
\sum_{n=1}^{99} \sqrt{1 + \frac{1}{n^2} + \frac{1}{(n+1)^2}}
= \sum_{n=1}^{99} \left(1 + \frac{1}{n} - \frac{1}{n+1}\right)
= 99 + 1 - \frac{1}{100}
= \frac{9999}{100}.
\]
\end{solution}

\begin{example}[$ \sum_{n=2}^{\infty} \left( \sqrt{ \frac{1}{F_{n-1}^2} + \frac{1}{F_n^2} + \frac{1}{F_{n+1}^2} } - \frac{1}{F_n} \right) $]
已知 \( F_n \) 为斐波那契数列。计算无穷级数：$$ \sum_{n=2}^{\infty} \left( \sqrt{ \frac{1}{F_{n-1}^2} + \frac{1}{F_n^2} + \frac{1}{F_{n+1}^2} } - \frac{1}{F_n} \right) $$
\end{example}
\begin{solution}
利用三项完全平方公式。当 \( x+y+z=0 \) 时，交叉项消失。斐波那契满足 \( F_{n-1} + F_n - F_{n+1} = 0 \)，令 \( x=F_{n-1}, y=F_n, z=-F_{n+1} \)。

由
\[
F_{n-1}+F_n-F_{n+1}=0
\]
可得
\[
\left|\frac{1}{F_{n-1}}+\frac{1}{F_n}-\frac{1}{F_{n+1}}\right|=\frac{1}{F_n},
\]
因此原项恒为 \(0\)，故原级数的和也为 \(0\)。
\end{solution}

\begin{example}
    $$ \sum_{n=1}^{\infty} \frac{1}{\sqrt{n^2 + (n+1)^2 + n^2(n+1)^2} - 1} $$
\end{example}
\begin{solution}
利用恒等式 \(n^2+(n+1)^2+n^2(n+1)^2=(n^2+n+1)^2\)，因此根号外再减去 \(1\) 后得到 \(n(n+1)\)。

原级数化为
\[
\sum_{n=1}^{\infty}\left(\frac{1}{n}-\frac{1}{n+1}\right)=1.
\]
\end{solution}

\begin{example}
    $$ \sum_{n=1}^{\infty} \frac{3}{\sqrt{n(n+1)(n+2)(n+3) + 1} - 1} $$
\end{example}
\begin{solution}
利用恒等式 \(n(n+1)(n+2)(n+3)+1=(n^2+3n+1)^2\)，因此根号减去 \(1\) 后得到 \(n(n+3)\)。

原级数化为
\[
\sum_{n=1}^{\infty} \left(\frac{1}{n}-\frac{1}{n+3}\right)
=1+\frac12+\frac13
=\frac{11}{6}.
\]
\end{solution}


\begin{example}[$\displaystyle \sum_{n=1}^{10} \sqrt{\frac{n^4 + (n+1)^4 + (2n+1)^4}{2}}$]
    $$\displaystyle \sum_{n=1}^{10} \sqrt{\frac{n^4 + (n+1)^4 + (2n+1)^4}{2}}$$
\end{example}
\begin{solution}
利用恒等式
\[
a^4+b^4+(a+b)^4=2(a^2+ab+b^2)^2.
\]
令 \(a=n,\ b=n+1\)，则根号内恰好化为完全平方，开方后得到
\[
3n^2+3n+1=(n+1)^3-n^3.
\]

根号化简为 \((n+1)^3 - n^3\)。求和得：
\[
\sum_{n=1}^{10} \bigl[(n+1)^3 - n^3\bigr] = 11^3 - 1^3 = 1331 - 1 = 1330
\]

\end{solution}

\begin{example}[$\displaystyle \sum_{n=1}^{99} \frac{1}{(n+1)\sqrt{n} + n\sqrt{n+1}}$]
    \[\displaystyle \sum_{n=1}^{99} \frac{1}{(n+1)\sqrt{n} + n\sqrt{n+1}}\]
\end{example}
\begin{solution}
从恒等式
\[
\frac{1}{\sqrt{n}}-\frac{1}{\sqrt{n+1}}
\]
出发通分，分母会出现题中的和式结构，于是原项被识别为相邻差。

原项等于 \(\frac{\sqrt{n+1} - \sqrt{n}}{\sqrt{n(n+1)}} \cdot \frac{\sqrt{n+1} + \sqrt{n}}{\sqrt{n+1} + \sqrt{n}} = \frac{1}{n} - \frac{1}{n+1}\)。求和得：
\[
1 - \frac{1}{100} = \frac{99}{100}
\]

\end{solution}

\begin{example}[$\displaystyle \sum_{n=1}^{99} \frac{n}{n^4+n^2+1}$]
\[\displaystyle \sum_{n=1}^{99} \frac{n}{n^4+n^2+1}\]
\end{example}
\begin{solution}
把 \(n^4 + n^2 + 1\) 写成 \((n^2+1)^2 - n^2 = (n^2-n+1)(n^2+n+1)\)。后一因式与下一项的前一因式相同，因为 \((n+1)^2 - (n+1) + 1 = n^2+n+1\)，于是设置相邻分式差：
\[
\frac{n}{n^4+n^2+1} = \frac{1}{2}\left(\frac{1}{n^2-n+1}-\frac{1}{n^2+n+1}\right).
\]

将部分分式展开后求和得：
\[
\frac{1}{2} \left(1 - \frac{1}{99^2 + 99 + 1}\right) = \frac{1}{2} \left(1 - \frac{1}{9901}\right) = \frac{4950}{9901}
\]

\end{solution}

\begin{example}[$\displaystyle \sum_{n=1}^{99} (n^2+n+1)n!$]
\[\displaystyle \sum_{n=1}^{99} (n^2+n+1)n!\]
\end{example}
\begin{solution}
令 \(f(n)=n\cdot n!\)。计算 \(f(n+1)-f(n)\) 并提取 \(n!\) 公因式，括号内得到 \(n^2+n+1\)。

将每一项写成 \(f(n+1)-f(n)\) 的形式后，求和得：
\[
100! - 1!
\]

\end{solution}

\section{递推型裂项}

\begin{example}
设数列满足 \(a_1=2, a_{n+1}=a_n^2-a_n+1\)。求无穷级数：
\[\displaystyle \sum_{n=1}^{\infty} \frac{1}{a_n}\]
\end{example}
\begin{solution}
利用恒等式
\[
\frac{1}{x-1}-\frac{1}{x}=\frac{1}{x(x-1)}.
\]
若令 \(a_{n+1}-1=a_n(a_n-1)\)，便得到题中的递推式 \(a_{n+1}=a_n^2-a_n+1\)。

原项可写为 \(\frac{1}{a_n-1} - \frac{1}{a_{n+1}-1}\)。求和得：
\[
\frac{1}{a_1-1} - \lim_{n\to\infty} \frac{1}{a_{n+1}-1} = 1 - 0 = 1
\]

\end{solution}

\begin{example}
   $$ \sum_{n=1}^{\infty} \frac{n^3-n^2-2n-1}{n!} $$
\end{example}
\begin{solution}
写下 \( \frac{n^3}{n!} - \frac{(n+1)^3}{(n+1)!} \)，上下同乘 \((n+1)\) 后约掉 \((n+1)\)，分子得到 \(n^3-n^2-2n-1\)。

原项等于 \( \frac{n^3}{n!} - \frac{(n+1)^3}{(n+1)!} \)。求和得：
\[
\frac{1^3}{1!} - \lim_{n\to\infty} \frac{(n+1)^3}{(n+1)!} = 1 - 0 = 1
\]

\end{solution}

\begin{example}
   设数列满足首项 \(a_1 = \sqrt{2}\)，且 \(a_{n+1} = a_n^3 - 3a_n\)。计算$$ \prod_{n=1}^{10} (a_n + 1) = (a_1+1)(a_2+1) \cdots (a_{10}+1) $$
\end{example}
\begin{solution}
递推式 \(a_{n+1}=a_n^3-3a_n\) 与三倍角公式相配，因为若 \(a_n=2\cos u\)，则
\[
a_{n+1}=2\cos 3u.
\]
取 \(u=2\cdot 3^{n-1}x\)，便有 \(a_n=2\cos(2\cdot3^{n-1}x)\)。再令 \(x=\pi/8\)，使 \(a_1=2\cos\frac{\pi}{4}=\sqrt2\)。由
\[
2\cos u+1=\frac{\sin\frac{3u}{2}}{\sin\frac{u}{2}},
\]
得到
\[
a_n+1=\frac{\sin(3^n x)}{\sin(3^{n-1}x)}.
\]

连乘后相邻项相消，得到
\[
\prod_{n=1}^{10} (a_n + 1) = \frac{\sin(3^{10} \cdot \pi/8)}{\sin(\pi/8)} = \frac{\sin(3^{10}\pi/8)}{\sin(\pi/8)}
\]
由于 \(3^{10}\equiv 9 \pmod{16}\)，故
\[
\sin\frac{3^{10}\pi}{8}=\sin\frac{9\pi}{8}=-\sin\frac{\pi}{8}.
\]
因此
\[
\prod_{n=1}^{10} (a_n + 1)=-1.
\]
\end{solution}

\section{指数型裂项}

本节处理一类指数型裂项。若把某一项写成
\[
\frac{1}{A_n}-\frac{k}{A_{n+1}},
\]
可先选取按 \(x\mapsto x^2\)、\(x\mapsto x^3\) 等方式迭代的基本量，使
\[
A_{n+1}=A_nQ_n.
\]
此时通分得到
\[
\frac{1}{A_n}-\frac{k}{A_nQ_n}=\frac{Q_n-k}{A_nQ_n}.
\]
若 \(Q_n-k\) 含有 \(A_n\) 因子，就得到裂项；常数 \(k\) 往往由令 \(A_n=0\) 确定。

例如在平方迭代中，若取 \(A=x-x^{-1}\)，则 \(A=0\) 对应 \(x=\pm1\)，这时 \(Q\) 取值不唯一。改取
\[
A=(x-x^{-1})^2
\]
后，\(A=0\) 等价于 \(x^2+x^{-2}=2\)。这时
\[
Q=(x+x^{-1})^2
\]
在该条件下唯一取值 \(4\)，于是得到
\[
\frac{1}{(x-x^{-1})^2}-\frac{4}{(x^2-x^{-2})^2}=\frac{1}{(x+x^{-1})^2}.
\]

在立方迭代 \(x\mapsto x^3\) 下，若取 \(A=x-1\)，则
\[
x^3-1=(x-1)(x^2+x+1).
\]
令 \(A=0\) 得 \(x=1\)，于是 \(k=3\)，从而
\[
\frac{1}{x-1}-\frac{3}{x^3-1}=\frac{x+2}{x^2+x+1}.
\]

更高次时也是同样的思路：先分解出上一层因子，再由常数项确定系数。


\begin{example}
    令 \(x_n = 2^{3^n}\)。计算无穷级数：$$ \sum_{n=1}^{\infty} 3^n \frac{x_n+2}{x_n^2+x_n+1} $$
\end{example}
\begin{solution}
利用立方递推 \(x\mapsto x^3\)。取 \(A=x-1\)，\(B=x^3-1=(x-1)(x^2+x+1)\)，\(Q=x^2+x+1\)。令 \(A=0\) 得 \(k=3\)，分子 \(Q-3\) 含 \(x-1\)，约去后得到 \(\frac{x+2}{x^2+x+1}\)。

原项等于 \(\frac{3^n}{x_n-1} - \frac{3^{n+1}}{x_{n+1}-1}\)。求和得：
\[
\frac{3}{8-1} - \lim_{n\to\infty} \frac{3^{n+1}}{2^{3^{n+1}}-1} = \frac{3}{7}
\]

\end{solution}

\begin{example}
    计算带有双重指数塔的无穷级数：$$ \sum_{n=0}^{\infty} 3^n \frac{4^{3^n}+2}{16^{3^n} + 4^{3^n} + 1} $$
\end{example}
\begin{solution}
沿用上一题的立方递推结构，可得 \(\frac{x^2+2}{x^4+x^2+1} = \frac{1}{x^2-1} - \frac{3}{x^6-1}\)。令 \(x=4^{3^n}\)。

因此原级数化为一列裂项和，其部分和满足
\[
\sum_{n=0}^{N} 3^n \frac{4^{3^n}+2}{16^{3^n}+4^{3^n}+1}
=\sum_{n=0}^{N}\left(\frac{3^n}{4^{2\cdot 3^n}-1}-\frac{3^{n+1}}{4^{2\cdot 3^{n+1}}-1}\right).
\]
令 \(N\to\infty\)，得原级数的和为
\[
\frac{1}{16-1}=\frac{1}{15}.
\]
\end{solution}

\begin{example}
    设数列满足 \(a_1 = 3\)，且 \(a_{n+1} = a_n^3 - 3a_n\)。计算无穷级数：$$ \sum_{n=1}^{\infty} 9^n \frac{a_n^2+2}{(a_n^2-1)^2} $$
\end{example}
\begin{solution}
令 \(A=y^2-4\)。由 \(y_{n+1}=y_n^3-3y_n\) 可推出
\[
y_{n+1}^2-4=(y_n^2-4)(y_n^2-1)^2.
\]
因此相邻两项通分时，分母自然出现 \((a_n^2-4)(a_n^2-1)^2\)。系数 \(9\) 使分子中最高阶部分约去，剩下的正是 \(a_n^2+2\)。

原级数化为
\[
\sum_{n=1}^{\infty}\left(\frac{9^n}{a_n^2-4}-\frac{9^{n+1}}{a_{n+1}^2-4}\right),
\]
故求和得
\[
\frac{9}{a_1^2-4}=\frac{9}{5}.
\]
\end{solution}

\begin{example}
    计算无穷连乘积：$$ \prod_{n=2}^{\infty} \frac{n^3-1}{n^3+1} = \frac{7}{9} \times \frac{26}{28} \times \frac{63}{65} \times \dots $$
\end{example}
\begin{solution}
展开 \(\frac{n^3-1}{n^3+1} = \frac{n-1}{n+1} \cdot \frac{n^2+n+1}{n^2-n+1}\)。定义 \(P(n)=n^2-n+1\)，则 \(P(n+1)=n^2+n+1\)。

将前 \(N\) 项连乘后，前一因子化为
\[
\prod_{n=2}^{N}\frac{n-1}{n+1}=\frac{2}{N(N+1)},
\]
后一因子化为
\[
\prod_{n=2}^{N}\frac{n^2+n+1}{n^2-n+1}=\frac{N^2+N+1}{3}.
\]
两部分相乘并令 \(N\to\infty\)，得极限为 \(\frac{2}{3}\)。
\end{solution}

\begin{example}
    已知常数 \(x\)，计算无穷级数：$$ \sum_{n=1}^\infty \frac{1}{2^n} \tan\left(\frac{x}{2^n}\right) = \frac{1}{2}\tan\left(\frac{x}{2}\right) + \frac{1}{4}\tan\left(\frac{x}{4}\right) + \dots $$
\end{example}
\begin{solution}
由二倍角公式变形 \(\cot A - 2\cot 2A = \tan A\)，再除以 \(2^n\) 得到指数权重裂项。

原级数化为
\[
\sum_{n=1}^{\infty}\left(\frac{1}{2^n}\cot\frac{x}{2^n}-\frac{1}{2^{n-1}}\cot\frac{x}{2^{n-1}}\right),
\]
故求和得 \(\frac{1}{x}-\cot x\)。
\end{solution}

\begin{example}
    已知非线性数列 \(y_0 = 3, y_{n+1} = y_n^3 - 3y_n\)。计算无穷级数：$$ \sum_{n=0}^\infty 9^n \frac{y_n^2+2}{(y_n^2-1)^2} $$
\end{example}
\begin{solution}
仍取 \(A=y_n^2-4\)。通分后分子可化为 \(A+6\)，从而得到标准裂项形式。

原级数化为
\[
\sum_{n=0}^{\infty}\left(\frac{9^n}{y_n^2-4}-\frac{9^{n+1}}{y_{n+1}^2-4}\right),
\]
故求和得 \(\frac{1}{5}\)。
\end{solution}

\begin{example}
设数列 \(y_0 = 3, y_{n+1} = y_n^4 - 4y_n^2 + 2\)。计算无穷级数：$$ \sum_{n=0}^\infty 16^n \frac{y_n^4+4}{y_n^2(y_n^2-2)^2} $$
\end{example}
\begin{solution}
取 \(A=y_n^2-4\)，构造 \(\frac{1}{A} - \frac{16}{A(A+4)(A+2)^2}\)，分子约去 \(A\) 后得到 \(\frac{A^2+8A+20}{(A+4)(A+2)^2}\)，代回即得 \(y_n^4+4\)。

原级数化为
\[
\sum_{n=0}^{\infty}\left(\frac{16^n}{y_n^2-4}-\frac{16^{n+1}}{y_{n+1}^2-4}\right),
\]
故求和得 \(\frac{1}{5}\)。
\end{solution}

\begin{example}
对于 \(0 < x < 1\)，计算无穷级数：$$ \sum_{n=0}^\infty \frac{4 x^{2^n} (x^{2^{n+1}} + x^{2^n} + 1)}{(1-x^{2^{n+1}})^2} $$
\end{example}
\begin{solution}
令
\[
f(u)=\frac{1+u}{1-u}.
\]
比较 \(f(u)\) 与 \(f(u^2)\)，可得
\[
f(u)-f(u^2)=\frac{2u}{1-u^2}.
\]
先计算
\[
f(u)^2-f(u^2)^2=\frac{4u(u^2+u+1)}{(1-u^2)^2}.
\]
再令 \(u=x^{2^n}\)，便得到题中的项。

原级数化为
\[
\sum_{n=0}^{\infty}\left[\left(\frac{1+x^{2^n}}{1-x^{2^n}}\right)^2-\left(\frac{1+x^{2^{n+1}}}{1-x^{2^{n+1}}}\right)^2\right],
\]
故求和得 \(\frac{4x}{(1-x)^2}\)。
\end{solution}

\begin{example}
已知常数 \(x \neq 0\)，求无穷级数：$$ \sum_{n=1}^{\infty} \frac{1}{4^n} \sec^2\left(\frac{x}{2^n}\right) $$
\end{example}
\begin{solution}
由 \(\sin 2\theta = 2\sin\theta\cos\theta\) 平方取倒数，并把分子写成 \(\sin^2\theta + \cos^2\theta\)，可得 \(\sec^2\theta = 4\csc^2 2\theta - \csc^2\theta\)。再除以 \(4^n\)。

原级数化为
\[
\sum_{n=1}^{\infty}\left(\frac{1}{4^{n-1}}\csc^2\frac{x}{2^{n-1}}-\frac{1}{4^n}\csc^2\frac{x}{2^n}\right),
\]
故求和得 \(\csc^2 x-\frac{1}{x^2}\)。
\end{solution}

三角函数中也常出现同样的裂项结构。由
\[
\tan A-\tan B=\frac{\sin(A-B)}{\cos A\cos B}
\]
在 \(A=(n+1)^\circ,\ B=n^\circ\) 时可得
\[
\frac{1}{\cos n^\circ\cos(n+1)^\circ}=\frac{1}{\sin1^\circ}\bigl(\tan(n+1)^\circ-\tan n^\circ\bigr).
\]
因此某些只含 \(\cos\) 分母的级数，也能写成正切差的裂项。

\begin{example}[$\displaystyle\sum_{n=1}^{88}\frac{1}{\cos n^\circ\cos(n+1)^\circ}$]
    计算求和：\[\sum_{n=1}^{88} \frac{1}{\cos n^\circ \cos (n+1)^\circ}\]
\end{example}
\begin{solution}
把 \(A=(n+1)^\circ,\ B=n^\circ\) 代入正切差公式，得到
\[
\tan(n+1)^\circ-\tan n^\circ=\frac{\sin1^\circ}{\cos n^\circ\cos(n+1)^\circ},
\]
分母就是相邻两个余弦的乘积，因此原项可写成相邻两项正切之差。

由上式可得，原项可写为
\[
\frac{1}{\sin 1^\circ}\bigl(\tan(n+1)^\circ-\tan n^\circ\bigr).
\]
求和得
\[
\frac{\tan 89^\circ - \tan 1^\circ}{\sin 1^\circ} = \frac{\cot 1^\circ - \tan 1^\circ}{\sin 1^\circ}.
\]

\end{solution}

连乘同样可这样处理。由
\[
\sin(2A)=2\sin A\cos A
\]
可得
\[
\cos A=\frac{\sin2A}{2\sin A},
\]
从而把连乘改写为相邻正弦项的连锁消去。

\begin{example}[$\cos20^\circ\cos40^\circ\cos80^\circ$]
    求纯数值表达式的值：\[\cos 20^\circ \cdot \cos 40^\circ \cdot \cos 80^\circ\]
\end{example}
\begin{solution}
利用
\[
\cos A=\frac{\sin2A}{2\sin A},
\]
把 \(\cos x\cdot\cos2x\cdot\cos4x\) 写成相邻正弦项的连乘分式。

在原式两侧补乘 \(8\sin 20^\circ\)，得
\[
\cos 20^\circ \cdot \cos 40^\circ \cdot \cos 80^\circ = \frac{\sin 160^\circ}{8\sin 20^\circ} = \frac{\sin 20^\circ}{8\sin 20^\circ} = \frac{1}{8}
\]

\end{solution}

类似的相邻角裂项也可用于三角函数。

\begin{example}[$\displaystyle\sum_{n=1}^{N} \frac{1}{\sin(2^n \theta)}$]
    设 \(N\in\mathbb N^*\)，且 \(\sin(2^n\theta)\ne0\ (1\le n\le N)\)，计算
    \[\sum_{n=1}^{N} \frac{1}{\sin(2^n \theta)}.\]
\end{example}
\begin{solution}
利用恒等式
\[
\cot x-\cot 2x=\frac{1}{\sin 2x},
\]
再把 \(x\) 替换为 \(2^{n-1}\theta\)，得到所需裂项。

原项可写为 \(\cot(2^{n-1}\theta) - \cot(2^n\theta)\)。有限求和时中间项全部消去，得到
\[
\sum_{n=1}^{N}\frac{1}{\sin(2^n\theta)}
=\cot\theta-\cot(2^N\theta).
\]

\end{solution}
